\chapter*{Conclusion générale}
\markboth{Conclusion générale}{}
\addcontentsline{toc}{chapter}{Conclusion générale}
\label{chap:05-conclusion}
Le projet a abouti à une application V1 couvrant l'enregistrement, la qualification, l'affectation, le traitement, la résolution et la clôture des demandes techniques. La séparation entre interface Angular, API Spring Boot et persistance PostgreSQL fournit une base lisible et évolutive. La sécurité combine authentification JWT, rôles métier et contrôles contextuels.

L'assistance IA facultative via Groq propose une catégorie et une priorité, mais la décision finale reste modifiable et validée par le Responsable technique. Les tests automatisés, les vérifications d'intégration réelle et les validations humaines confirment le fonctionnement du périmètre V1 dans les conditions vérifiées.

Ce stage a permis de mettre en pratique une démarche complète de génie logiciel, depuis la compréhension du besoin et sa modélisation UML jusqu'à la conception, la réalisation, les tests et la validation. Il a également renforcé la distinction entre les besoins fonctionnels et les choix techniques, ainsi que la nécessité de justifier une architecture par les responsabilités qu'elle doit séparer. La réalisation individuelle du projet a demandé de l'autonomie et une progression contrôlée, notamment pour articuler Angular, Spring Boot et PostgreSQL, mettre en œuvre les contrôles JWT et RBAC, puis intégrer un service IA externe sans lui confier la décision métier.

Les limites concernent le périmètre interne, l'absence de portail client et le caractère exploratoire du benchmark IA. Des travaux ultérieurs pourraient élargir le corpus d'évaluation, renforcer l'observabilité et étudier les besoins de déploiement, sous réserve d'une analyse du besoin et de validations supplémentaires.
